\documentclass[10pt,a4paper]{article}

% Solucionario del reto de la situación de aprendizaje «Medir lo que no se puede alcanzar».
% Reutiliza el preámbulo y la bibliografía de la situación de aprendizaje.
% Compilar desde esta carpeta con: BIBINPUTS=..: latexmk -pdf Solucionario_SA_JGL_DCM.tex

% ---------- Datos del documento ----------
\newcommand{\tituloSA}{Medir lo que no se puede alcanzar: ¿dónde caben más paneles solares en nuestro instituto?}
\newcommand{\autor}{Javier Gómez López}
\newcommand{\asignatura}{Diseño Curricular de las Matemáticas}
\newcommand{\master}{Máster Universitario en Formación del Profesorado -- Especialidad en Matemáticas}

% ---------- Paquetes ----------
\usepackage{float}
\usepackage[utf8]{inputenc}
\usepackage[T1]{fontenc}
\usepackage[spanish,es-noshorthands,es-tabla]{babel}
\usepackage[scaled=0.92]{helvet} % equivalente libre de la Arial del original
\renewcommand{\familydefault}{\sfdefault}
\usepackage[left=2cm,right=2cm,top=2.55cm,bottom=1.6cm,
            headheight=44pt,headsep=4pt,footskip=20pt]{geometry}
\usepackage{amsmath,amssymb}
\usepackage{xcolor}
\usepackage{graphicx}
\usepackage{fancyhdr}
\usepackage{lastpage}
\usepackage{etoolbox}
\usepackage{enumitem}
\usepackage{titlesec}
\usepackage{pdflscape}
\usepackage{tikz}
\usetikzlibrary{angles,quotes,arrows.meta}
\usepackage{tabularray}
\usepackage[most]{tcolorbox}
\usepackage{csquotes}
\usepackage[backend=bibtex,style=authoryear,maxcitenames=2,uniquename=false,uniquelist=false]{biblatex}
\addbibresource{bibliografia.bib}
\usepackage[hidelinks]{hyperref}

\setlength{\parindent}{0pt}
\setlength{\parskip}{5pt}
\setlist{nosep,leftmargin=1.2em}
\emergencystretch=1.5em

% ---------- Colores (medidos sobre la plantilla) ----------
\definecolor{urjc}{HTML}{C8102E}
\definecolor{rosa}{HTML}{FBE9EB}
\definecolor{borde}{HTML}{D0D0D0}
\definecolor{titular}{HTML}{3C3C3B}
\definecolor{pistagris}{HTML}{7F7F7F}

% ---------- Títulos ----------
\titleformat{\section}{\Large\bfseries\color{urjc}}{\thesection}{0.6em}{}[{\color{urjc}\titlerule[0.7pt]}]
\titleformat{\subsection}{\large\bfseries\color{titular}}{\thesubsection}{0.6em}{}
\titleformat{\subsubsection}{\normalsize\bfseries\color{urjc}}{}{0pt}{}
\titlespacing*{\section}{0pt}{14pt}{6pt}
\titlespacing*{\subsection}{0pt}{10pt}{3pt}
\titlespacing*{\subsubsection}{0pt}{8pt}{2pt}
\setcounter{secnumdepth}{2}
\setcounter{tocdepth}{2}
\makeatletter\renewcommand*\l@subsection{\@dottedtocline{2}{1.5em}{2.9em}}\makeatother

% ---------- Cabecera y pie ----------
\newcommand{\cabecera}{Máster en Formación del Profesorado · Diseño Curricular de las Matemáticas · 2026-27}

\pagestyle{fancy}
\fancyhf{}
\fancyhead[L]{\includegraphics[height=42pt]{logo_urjc.png}}
\fancyhead[R]{\footnotesize\color{pistagris}\cabecera}
\fancyfoot[R]{\footnotesize\color{pistagris}Página \thepage\ de \pageref{LastPage}}
\renewcommand{\headrule}{{\color{urjc}\hrule height 1pt}}
\renewcommand{\footrulewidth}{0pt}

% ---------- Elementos de texto ----------
% Etiqueta de un campo (Título:, Producto:...)
\newcommand{\campo}[1]{{\bfseries\color{urjc}#1}}
% Texto orientativo, en cursiva gris
\newcommand{\pista}[1]{{\itshape\small\color{pistagris}#1}}
% Lista con viñetas; dentro de las tablas hay que usar esta y no itemize
\newenvironment{lista}
  {\begin{minipage}[t]{\linewidth}\begin{itemize}[leftmargin=*,nosep]}
  {\end{itemize}\end{minipage}}
% Encabezado de los apartados a todo el ancho de la plantilla (7 y 8)
\newcommand{\apartado}[1]{%
  \par\vspace{14pt}%
  {\large\bfseries\color{urjc}#1\par}%
  \vspace{4pt}{\color{urjc}\hrule height 0.7pt}\vspace{5pt}}

% ---------- Recuadros ----------
% \begin{tarjeta}[subtítulo opcional]{título}
\newtcolorbox{tarjeta}[2][]{%
  enhanced, sharp corners, boxrule=0.5pt, titlerule=0pt,
  colframe=borde, colback=white, colbacktitle=urjc, coltitle=white,
  fonttitle=\bfseries\large, toptitle=3pt, bottomtitle=3pt,
  left=7pt, right=7pt, top=7pt, bottom=7pt,
  halign=flush left, halign title=flush left,
  before upper={\setlength{\parskip}{5pt}},
  title={#2\ifstrempty{#1}{}{\\[-1pt]{\footnotesize\itshape #1}}}}

% Fila de dos recuadros de igual altura
\newenvironment{filatarjetas}{%
  \begin{tcbraster}[raster columns=2, raster equal height=rows,
                    raster column skip=8pt, raster after skip=12pt]}
  {\end{tcbraster}}

% Sesión de la secuencia: \sesion{n}{tipo}{título}{duraciones y texto}
\newtcolorbox{sesionbox}[3]{%
  enhanced, sharp corners, boxrule=0.5pt, titlerule=0pt,
  colframe=borde, colback=white, colbacktitle=urjc, coltitle=white,
  fonttitle=\bfseries, toptitle=2pt, bottomtitle=2pt,
  left=6pt, right=6pt, top=4pt, bottom=5pt, after skip=8pt,
  halign=flush left, halign title=flush left,
  fontupper=\small, before upper={\setlength{\parskip}{3pt}},
  title={Sesión #1 · #2: #3}}

% Caja de información
\newtcolorbox{nota}[1]{%
  enhanced, sharp corners, boxrule=0pt, leftrule=3pt, colframe=urjc, colback=rosa,
  coltitle=urjc, colbacktitle=rosa, fonttitle=\bfseries, attach title to upper={\ },
  left=7pt, right=7pt, top=4pt, bottom=4pt, title={#1}, fontupper=\small}

% ---------- Tablas ----------
\NewTblrEnviron{tablaSA}
\SetTblrInner[tablaSA]{
  hlines={0.5pt,borde}, vlines={0.5pt,borde},
  leftsep=5pt, rightsep=5pt, rowsep=3pt,
  cells={font=\footnotesize, valign=t},
  row{1}={bg=rosa, fg=urjc, font=\footnotesize\bfseries},
}

% Versión con salto de página y cabecera repetida
\NewTblrTheme{limpio}{
  \SetTblrTemplate{caption}{empty}
  \SetTblrTemplate{capcont}{empty}
  \SetTblrTemplate{contfoot}{empty}
  \SetTblrTemplate{conthead}{empty}
  \SetTblrTemplate{note}{empty}
  \SetTblrTemplate{remark}{empty}
}
\NewTblrEnviron{tablaSAlarga}
\SetTblrOuter[tablaSAlarga]{long, theme=limpio, entry=none, label=none, presep=2pt, postsep=4pt}
\SetTblrInner[tablaSAlarga]{
  hlines={0.5pt,borde}, vlines={0.5pt,borde},
  leftsep=5pt, rightsep=5pt, rowsep=3pt,
  cells={font=\footnotesize, valign=t},
  rowhead=1,
}

\graphicspath{{../}}
\renewcommand{\cabecera}{Solucionario del reto · Diseño Curricular de las Matemáticas · 2026-27}

% Resultados y demostraciones
\newcounter{res}
\newcommand{\resultado}[2]{\par\medskip\noindent\refstepcounter{res}{\bfseries\color{urjc}Resultado \theres. #1.} \emph{#2}\par\smallskip}
\newenvironment{dem}{\noindent{\itshape Demostración.} }{\hfill$\square$\par\medskip}
\newcommand{\paso}[1]{\par\medskip\noindent{\bfseries\color{titular}#1}\par\smallskip}

\begin{document}

% ---------- Portada ----------
\thispagestyle{empty}
\vspace*{3.2cm}
\begin{center}
  {\small\color{pistagris}\MakeUppercase{\master}\\ Universidad Rey Juan Carlos}\\[1.4em]
  {\color{urjc}\rule{\linewidth}{1.2pt}}\\[1.2em]
  {\fontsize{13}{16}\selectfont\color{titular} Solucionario del reto y demostraciones}\\[0.8em]
  {\fontsize{25}{31}\selectfont\bfseries\color{titular}\tituloSA\par}
  \vspace{1em}
  {\color{urjc}\rule{\linewidth}{1.2pt}}\\[1.4em]
  {\large Matemáticas I · 1.º de Bachillerato (Ciencias y Tecnología)}\\[0.4em]
  {\large\color{pistagris}\asignatura}\\[3em]
  {\large\bfseries\autor}\\[0.4em]
  {\small\color{pistagris}Curso 2026-27}
\end{center}
\clearpage

{\setlength{\parskip}{0pt}\tableofcontents}
\clearpage

% ======================================================================
\section*{Presentación}
\addcontentsline{toc}{section}{Presentación}

Este documento acompaña a la situación de aprendizaje y está pensado para el docente. Tiene dos partes:

\begin{itemize}
  \item \textbf{Guía paso a paso (secciones 1 a~\ref{sec:cierre}).} Resuelve el reto completo tal como lo haría un equipo: qué medir, cómo calcular, cómo construir el modelo en GeoGebra y cómo llegar a la propuesta. Sigue el orden de las partes del libro de GeoGebra y de las sesiones de trabajo (S13 a S17).
  \item \textbf{Demostraciones (sección~\ref{sec:demos}).} Justifica cada fórmula que usa el alumnado, desde la definición de las razones trigonométricas hasta el criterio de separación entre filas.
\end{itemize}

\begin{nota}{Datos de ejemplo.}
Las medidas, el plano y las dimensiones de la cubierta son inventados, pero coherentes entre sí y con la situación de aprendizaje, y todos los cálculos están comprobados. Cada equipo obtendrá valores distintos según su centro y sus medidas. Las soluciones de las fichas F1 a F4, del cuestionario y del examen ya figuran en la propia situación de aprendizaje.
\end{nota}

\paso{Notación común}
\begin{tablaSA}{colspec={Q[l,wd=50pt]X[l]}, cells={font=\footnotesize}}
  Símbolo & Significado \\
  $a$ & Altura del ojo de quien mide (m) \\
  $D$, $d$ & Distancia al pie del objeto; distancia entre dos puntos de medida (m) \\
  $\alpha$, $\beta$ & Ángulos de elevación medidos con el clinómetro, con $\alpha<\beta$ \\
  $H$ & Altura del objeto (m) \\
  $\varphi$, $\delta$ & Latitud del centro y declinación solar del día \\
  $h$, $A$ & Altura del sol sobre el horizonte y ángulo que forma con el sur \\
  $L$, $k$ & Longitud de la sombra y coeficiente $k=L/H=1/\tan h$ \\
  $\beta_p$, $\theta$ & Inclinación del panel y ángulo de incidencia del sol sobre él \\
\end{tablaSA}

% ======================================================================
\part*{Guía paso a paso}
\addcontentsline{toc}{part}{Guía paso a paso}

\section{Preparar la medición (S13)}

Antes de salir al patio, cada equipo decide qué objetos mide y con qué método, y el profesor lo aprueba.

\paso{Paso 1. Elegir el objeto y el método}
\begin{itemize}
  \item \textbf{Si se puede llegar al pie del objeto}, como en el edificio principal, basta \emph{un punto de medida}: se mide la distancia $D$ con la cinta y el ángulo de elevación $\alpha$.
  \item \textbf{Si no se puede llegar}, como en el bloque de viviendas tras la valla, hacen falta \emph{dos puntos de medida} $A$ y $B$, alineados con el objeto y separados una distancia $d$ que sí se puede medir.
  \item Además, se mide el edificio principal también con dos puntos, para comparar ambos métodos en un objeto cuya altura ya se conoce por el otro.
\end{itemize}

\paso{Paso 2. Calibrar el clinómetro}
Se apunta a un objeto horizontal, como el borde de una mesa a la altura del ojo: el clinómetro debe marcar $0^\circ$. Si el clinómetro casero mide el ángulo con la vertical (transportador con plomada), el ángulo de elevación es $\alpha = 90^\circ - \text{lectura}$.

\begin{nota}{Errores frecuentes.}
Olvidar sumar la altura del ojo $a$; usar la calculadora en radianes; medir $d$ en una dirección que no está alineada con el objeto; confundir qué punto es $A$ (el más lejano, con el ángulo menor $\alpha$) y cuál es $B$.
\end{nota}

% ======================================================================
\section{Medir el instituto (S14)}

\paso{Paso 3. Tomar y registrar los datos}
Cada magnitud se mide tres veces y se elige la media. Hoja de datos completa del equipo de ejemplo:

\begin{tablaSA}{colspec={X[2.6,l]X[1,c]X[1,c]X[1,c]X[1,c]}, cells={font=\footnotesize}}
  Magnitud (con unidad) & Medida 1 & Medida 2 & Medida 3 & Valor elegido \\
  Altura del ojo $a$ (m) & 1,60 & 1,61 & 1,59 & 1,60 \\
  \SetCell[c=5]{l}{\bfseries Edificio principal, un punto de medida} & & & & \\
  Distancia al pie $D$ (m) & 24,78 & 24,82 & 24,80 & 24,80 \\
  Ángulo de elevación $\alpha$ (°) & 25,5 & 24,5 & 25,0 & 25,0 \\
  \SetCell[c=5]{l}{\bfseries Edificio principal, dos puntos de medida ($B$ es el punto anterior)} & & & & \\
  Distancia entre puntos $d$ (m) & 10,00 & 10,01 & 9,99 & 10,00 \\
  Ángulo $\alpha$ en $A$ (°) & 18,5 & 18,2 & 18,5 & 18,4 \\
  Ángulo $\beta$ en $B$ (°) & 25,5 & 24,5 & 25,0 & 25,0 \\
  \SetCell[c=5]{l}{\bfseries Bloque de viviendas, dos puntos de medida} & & & & \\
  Distancia entre puntos $d$ (m) & 10,00 & 10,02 & 9,98 & 10,00 \\
  Ángulo $\alpha$ en $A$ (°) & 24,5 & 25,5 & 25,0 & 25,0 \\
  Ángulo $\beta$ en $B$ (°) & 38,5 & 37,5 & 38,0 & 38,0 \\
  Sombra de la varilla de 1~m (m), 10 de marzo, mediodía solar & 1,00 & 1,02 & 1,01 & 1,01 \\
\end{tablaSA}

\paso{Paso 4. Validar los datos antes de irse}
Se revisan las tres medidas de cada fila: si una se aleja mucho de las otras dos (por ejemplo, más de un grado), se repite. En el ejemplo todas las filas son coherentes.

% ======================================================================
\section{Calcular las alturas (S15)}

\begin{figure}[H]
\centering
\begin{minipage}[b]{0.49\linewidth}\centering
\begin{tikzpicture}[>=Stealth,font=\footnotesize]
  \fill[gray!30] (4.6,0) rectangle (5.4,2.745); \draw[thick] (4.6,0) rectangle (5.4,2.745);
  \draw[thick] (-0.5,0) -- (5.9,0);
  \draw[thick] (0,0) -- (0,0.6); \fill (0,0.6) circle (1.5pt);
  \draw[dashed] (0,0.6) -- (4.6,0.6);
  \draw[urjc,thick] (0,0.6) -- (4.6,2.745); \fill[urjc] (4.6,2.745) circle (1.5pt) node[above left]{$T$};
  \draw (1.1,0.6) arc (0:25:1.1); \node at (1.5,0.87) {$\alpha$};
  \draw[<->] (-0.3,0) -- node[left]{$a$} (-0.3,0.6);
  \draw[<->] (0,-0.3) -- node[below]{$D$} (4.6,-0.3);
  \draw[<->] (4.35,0.6) -- node[left]{$x$} (4.35,2.745);
  \draw[<->] (5.65,0) -- node[right]{$H$} (5.65,2.745);
\end{tikzpicture}
\par\smallskip (a)
\end{minipage}\hfill
\begin{minipage}[b]{0.49\linewidth}\centering
\begin{tikzpicture}[>=Stealth,font=\footnotesize]
  \fill[gray!30] (4.962,0) rectangle (5.7,2.634); \draw[thick] (4.962,0) rectangle (5.7,2.634);
  \fill[urjc!12] (0,0.32) -- (2,0.32) -- (4.962,2.634) -- cycle;
  \draw[thick] (-0.4,0) -- (6.1,0);
  \draw[thick] (0,0) -- (0,0.32) (2,0) -- (2,0.32);
  \fill (0,0.32) circle (1.5pt) node[above left]{$A$}; \fill (2,0.32) circle (1.5pt) node[above left]{$B$};
  \draw[dashed] (0,0.32) -- (4.962,0.32); \node[below right] at (4.962,0.32) {$F$};
  \draw[urjc,thick] (0,0.32) -- (4.962,2.634) (2,0.32) -- (4.962,2.634);
  \fill[urjc] (4.962,2.634) circle (1.5pt) node[above left]{$T$};
  \draw (0.95,0.32) arc (0:25:0.95); \node at (1.3,0.5) {$\alpha$};
  \draw (2.6,0.32) arc (0:38:0.6); \node at (2.88,0.58) {$\beta$};
  \draw[<->] (0,-0.3) -- node[below]{$d$} (2,-0.3);
\end{tikzpicture}
\par\smallskip (b)
\end{minipage}
\caption{Altura con un punto de medida (a) y con dos puntos de medida (b). $T$ es lo alto del objeto y $F$, el punto de su vertical a la altura del ojo.}\label{fig:alturas}
\end{figure}

\paso{Paso 5. Edificio principal con un punto de medida}
En el triángulo rectángulo de la figura~\ref{fig:alturas}a, $D$ es el cateto adyacente a $\alpha$ y $x$ el opuesto, así que $x=D\tan\alpha$ (resultado~\ref{res:unpunto}):
\[
  x = 24{,}80\cdot\tan 25^\circ \approx 11{,}56\ \text{m}, \qquad H = x + a \approx 11{,}56 + 1{,}60 = 13{,}16\ \text{m}.
\]

\paso{Paso 6. Bloque de viviendas con dos puntos de medida}
No se conoce la distancia al pie, pero sí $d$. Se trabaja en el triángulo $ABT$ (figura~\ref{fig:alturas}b, resultado~\ref{res:dospuntos}):
\begin{enumerate}[leftmargin=1.6em]
  \item Ángulos del triángulo: en $A$, $\alpha=25^\circ$; en $B$, $180^\circ-\beta=142^\circ$; en $T$, $\beta-\alpha=13^\circ$.
  \item Teorema del seno: $\dfrac{BT}{\sin\alpha}=\dfrac{d}{\sin(\beta-\alpha)}$, luego $BT=\dfrac{10\cdot\sin 25^\circ}{\sin 13^\circ}\approx 18{,}79$~m.
  \item Triángulo rectángulo $BFT$: $x = BT\sin\beta \approx 18{,}79\cdot\sin 38^\circ \approx 11{,}57$~m.
  \item Altura: $H = x + a \approx 11{,}57+1{,}60 = 13{,}17$~m.
\end{enumerate}
Todo junto: $H=\dfrac{d\,\sin\alpha\,\sin\beta}{\sin(\beta-\alpha)}+a$. Un camino equivalente, con dos triángulos rectángulos y tangentes, da $x=\dfrac{d\tan\alpha\tan\beta}{\tan\beta-\tan\alpha}\approx 11{,}57$~m, el mismo resultado.

\paso{Paso 7. Comparar los dos métodos en el edificio principal}
Con $d=10$~m, $\alpha=18{,}4^\circ$ y $\beta=25{,}0^\circ$:
\[
  H=\frac{10\cdot\sin 18{,}4^\circ\cdot\sin 25^\circ}{\sin 6{,}6^\circ}+1{,}60\approx 13{,}21\ \text{m},
\]
frente a los $13{,}16$~m del método de un punto. Los dos métodos coinciden en unos 5~cm, lo que valida el procedimiento de dos puntos antes de aplicarlo al bloque, donde no hay otra forma de comprobarlo.

\paso{Paso 8. Validar el modelo del sol con la varilla}
El 10 de marzo es el día $d=69$ del año. La declinación se obtiene con la fórmula de PVEducation \parencite{pveducationdecl}:
\[
  \delta=-23{,}45^\circ\cos\!\left[\frac{360^\circ}{365}(d+10)\right]\approx -4{,}91^\circ,\qquad h=90^\circ-\varphi+\delta\approx 90^\circ-40{,}4^\circ-4{,}91^\circ=44{,}69^\circ.
\]
El modelo predice una sombra de $1/\tan 44{,}69^\circ\approx 1{,}011$~m para la varilla de 1~m. La medida fue $1{,}01$~m, que corresponde a $h=\arctan(1/1{,}01)\approx 44{,}71^\circ$. Modelo y realidad difieren en dos centésimas de grado: el modelo funciona.

% ======================================================================
\section{Analizar el error (S15)}

\paso{Paso 9. Efecto de un grado de error}
Se repite cada cálculo cambiando un ángulo en $\pm 1^\circ$:

\begin{tablaSA}{colspec={X[2.2,l]X[1.4,c]X[1,c]X[1,c]X[1,c]}, cells={font=\footnotesize}}
  Objeto y método & Ángulo alterado & $-1^\circ$ & Valor medido & $+1^\circ$ \\
  Edificio, un punto & $\alpha$ & 12,64~m & 13,16~m & 13,70~m \\
  Edificio, dos puntos & $\alpha$ & 15,99~m & 13,21~m & 11,16~m \\
  Bloque, dos puntos & $\alpha$ & 11,95~m & 13,17~m & 14,58~m \\
  Bloque, dos puntos & $\beta$ & 13,83~m & 13,17~m & 12,59~m \\
\end{tablaSA}

\paso{Paso 10. Interpretar y elegir el resultado}
\begin{itemize}
  \item Con un punto de medida, un grado de error cambia la altura en algo más de medio metro. Con dos, en el bloque cambia más de un metro, y en el edificio casi tres, porque allí $\beta-\alpha$ es solo de $6{,}6^\circ$. El motivo es el denominador $\sin(\beta-\alpha)$: cuanto más parecidos son los dos ángulos, más pequeño es y más amplifica el error (resultado~\ref{res:error}).
  \item \textbf{Conclusión del equipo:} se usa un punto de medida siempre que se pueda; con dos, conviene separar bastante $A$ y $B$ para que $\beta-\alpha$ sea grande.
  \item \textbf{Resultados elegidos:} edificio principal, $13{,}2$~m; bloque de viviendas, $13{,}2$~m, sabiendo que podría estar entre unos $12$ y $14{,}6$~m. En la propuesta se comprueba si la decisión cambia dentro de ese intervalo.
\end{itemize}

% ======================================================================
\section{Construir el modelo en GeoGebra (S15)}

\paso{Paso 11. La construcción}
Se escribe en la barra de entrada, en este orden. Los nombres de los comandos son los de GeoGebra en español; en inglés son \texttt{Slider}, \texttt{Ray}, \texttt{Intersect} y \texttt{xAxis}.

\begin{tablaSA}{colspec={Q[l,wd=170pt]X[l]}, cells={font=\footnotesize}}
  Entrada & Qué hace \\
  \texttt{lat = 40.4°} & Latitud del centro, $\varphi$ \\
  \texttt{d = Deslizador(1, 365, 1)} & Día del año: deslizador para la fecha \\
  \texttt{dec = -23.45° cos(360° / 365 (d + 10))} & Declinación del día, $\delta$ \\
  \texttt{h = 90° - lat + dec} & Altura del sol al mediodía \\
  \texttt{H = Deslizador(0, 30, 0.1)} & Altura del obstáculo, reutilizable \\
  \texttt{T = (0, H)} & Lo alto del obstáculo \\
  \texttt{v = (cos(h), -sin(h))} & Vector director del rayo \\
  \texttt{r = Semirrecta(T, T + v)} & Rayo de sol que roza el obstáculo \\
  \texttt{S = Interseca(r, EjeX)} & Final de la sombra \\
  \texttt{L = x(S)} & Longitud de la sombra \\
  \texttt{k = 1 / tan(61° - lat)} & Coeficiente del criterio del IDAE \\
  \texttt{dIDAE = k H} & Separación mínima según el IDAE \\
\end{tablaSA}

\paso{Paso 12. Comprobación a mano}
Con $H=12$ se mueve el deslizador de la fecha y se compara con el cálculo $L=12/\tan h$:

\begin{tablaSA}{colspec={X[1.6,l]X[1,c]X[1,c]X[1,c]X[1,c]}, cells={font=\footnotesize}}
  Fecha & Día $d$ & $\delta$ & $h$ & $L$ para $H=12$~m \\
  21 de diciembre & 355 & $-23{,}45^\circ$ & $26{,}15^\circ$ & 24,44~m \\
  22 de marzo & 81 & $-0{,}10^\circ$ & $49{,}50^\circ$ & 10,25~m \\
  21 de junio & 172 & $+23{,}45^\circ$ & $73{,}05^\circ$ & 3,66~m \\
\end{tablaSA}

En el equinoccio la fórmula da $\delta\approx-0{,}10^\circ$ en lugar de $0^\circ$, porque es una aproximación: la diferencia en la sombra es de unos 4~cm.

\paso{Paso 13. Prueba de reutilización}
Se cambian la fecha y la altura y se comprueba que el modelo responde, por ejemplo con el bloque ($H=13{,}2$) el 21 de diciembre: $L\approx 26{,}8$~m y $\text{dIDAE}\approx 35{,}0$~m. Así la dirección puede usar el archivo con cualquier obstáculo.

% ======================================================================
\section{Elaborar la propuesta (S16)}

El equipo estudia dos zonas: la franja del patio junto a la valla sur y la cubierta plana del edificio principal (figura~\ref{fig:plano}). Las separaciones se calculan con el criterio del IDAE, $k\approx 2{,}66$, y se comparan con las del mediodía del solsticio, $k\approx 2{,}04$.

\begin{figure}[!htb]
\centering
\begin{tikzpicture}[x=0.19cm,y=0.115cm,>=Stealth,font=\footnotesize]
  % de sur (abajo) a norte (arriba), en metros
  \fill[gray!35] (0,0) rectangle (40,4); \node at (20,2) {bloque de viviendas ($H\approx 13{,}2$ m)};
  \draw[thick,dashed] (-2,10) -- (42,10); \node[right] at (42,10) {valla};
  \fill[green!12] (5,10) rectangle (35,30); \draw (5,10) rectangle (35,30);
  \node[align=center] at (20,20) {franja del patio\\$30\times 20$ m};
  \fill[gray!15] (0,30) rectangle (40,46); \node at (20,38) {patio de deportes};
  \fill[gray!35] (0,46) rectangle (40,62); \draw[very thick] (0,46) rectangle (40,62);
  \node[align=center] at (20,54) {cubierta del edificio principal\\$40\times 16$ m, $H\approx 13{,}2$ m};
  \draw[<->] (-3,4) -- node[left]{6 m} (-3,10);
  \draw[<->] (-3,4) -- (-3,46); \node[left] at (-3,28) {42 m};
  \draw[->,thick] (46,50) -- (46,58) node[above]{N};
\end{tikzpicture}
\caption{Plano esquemático del centro de ejemplo, de sur (abajo) a norte (arriba). Las distancias se miden desde la fachada norte del bloque.}\label{fig:plano}
\end{figure}

\paso{Paso 14. Primera zona: la franja del patio}
El bloque está al sur de la franja, a 6~m de su borde. Los paneles irían sobre el suelo, así que el obstáculo tiene toda su altura:
\begin{itemize}
  \item Criterio del IDAE: $2{,}66\cdot 13{,}17\approx 35{,}0$~m desde el bloque, es decir, la primera fila empezaría a $29{,}0$~m del borde de la franja, que solo tiene 20~m de fondo.
  \item Incluso con el sol del mediodía haría falta $2{,}04\cdot 13{,}17\approx 26{,}8$~m, y la fila empezaría a $20{,}8$~m: tampoco cabe.
  \item Con la altura más baja del intervalo de error, $11{,}95$~m, la primera fila seguiría a $25{,}8$~m del borde.
\end{itemize}
\textbf{Conclusión:} la franja no es viable por la sombra del bloque, y la decisión no cambia por el error de medida. Es justo el resultado que exigía medir una altura inaccesible.

\paso{Paso 15. Segunda zona: la cubierta del edificio principal}
\begin{enumerate}[leftmargin=1.6em]
  \item \textbf{Sombra del bloque.} El bloque ($13{,}17$~m) y el edificio ($13{,}16$~m) tienen casi la misma altura, así que el bloque no proyecta sombra sobre la cubierta. En el peor caso del error, el bloque mediría $14{,}58$~m y sobresaldría $1{,}42$~m: su sombra útil llegaría a $2{,}66\cdot 1{,}42\approx 3{,}8$~m, muy lejos de los 42~m que lo separan del edificio.
  \item \textbf{Petril.} La cubierta tiene un petril de 1~m en su borde sur. La primera fila debe quedar a $2{,}66\cdot 1\approx 2{,}66$~m de él.
  \item \textbf{Geometría de cada fila.} Paneles de $2{,}28\times 1{,}13$~m (valores ilustrativos), colocados en horizontal con el lado de $1{,}13$~m en la pendiente e inclinados $\beta_p=30^\circ$. Cada fila alcanza una altura $h_o=1{,}13\sin 30^\circ=0{,}565$~m y ocupa en planta $1{,}13\cos 30^\circ\approx 0{,}98$~m.
  \item \textbf{Paso entre filas} (resultado~\ref{res:filas}): $p=0{,}98+2{,}66\cdot 0{,}565\approx 2{,}48$~m.
  \item \textbf{Número de filas.} Se deja 1~m libre en el borde norte, así que hay sitio entre $2{,}66$~m y $15$~m desde el petril. Caben 5 filas: la última termina a $2{,}66+4\cdot 2{,}48+0{,}98\approx 13{,}57$~m.
  \item \textbf{Paneles por fila.} Con 1~m libre a cada lado quedan 38~m: $38/2{,}28\approx 16{,}7$, es decir, 16 paneles.
  \item \textbf{Total:} $5\cdot 16=80$ paneles. Con paneles de 550~W (valor ilustrativo), $80\cdot 0{,}55=44$~kWp.
\end{enumerate}

\begin{figure}[!htb]
\centering
\begin{tikzpicture}[x=1.45cm,y=1.45cm,>=Stealth,font=\footnotesize]
  \draw[thick] (-0.5,0) -- (7.0,0);
  \fill[gray!35] (-0.2,0) rectangle (0,1); \draw[thick] (-0.2,0) rectangle (0,1);
  \node[left] at (-0.2,0.5) {petril};
  \draw[blue!60!black,very thick] (2.6605,0) -- (3.6391,0.565);
  \draw[blue!60!black,very thick] (5.1423,0) -- (6.1209,0.565);
  \draw[urjc,thick] (0,1) -- (2.6605,0);
  \draw[urjc,thick] (3.6391,0.565) -- (5.1423,0);
  \draw (4.5423,0) arc (180:159.4:0.6); \node at (4.25,0.13) {$20{,}6^\circ$};
  \draw[<->] (0,-0.25) -- node[below]{$k\cdot 1\approx 2{,}66$ m} (2.6605,-0.25);
  \draw[<->] (2.6605,-0.75) -- node[below]{paso $p\approx 2{,}48$ m} (5.1423,-0.75);
  \draw[<->] (3.6391,-0.25) -- node[below]{$k\,h_o\approx 1{,}50$ m} (5.1423,-0.25);
  \draw[dashed] (3.6391,0.565) -- (3.6391,-0.3);
  \draw[<->] (6.4,0) -- node[right]{$h_o=0{,}565$ m} (6.4,0.565);
  \draw (3.2605,0) arc (0:30:0.6); \node at (3.36,0.09) {$30^\circ$};
  \node[below] at (-0.3,-0.25) {sur}; \node[below] at (6.8,-0.25) {norte};
\end{tikzpicture}
\caption{Corte sur-norte de la cubierta con el petril y las dos primeras filas. El sol de diseño del IDAE ($20{,}6^\circ$) roza lo alto de cada obstáculo y llega justo al pie de la fila siguiente.}\label{fig:filas}
\end{figure}

\paso{Paso 16. Comparar con el criterio del mediodía}
Con $k\approx 2{,}04$ la primera fila podría estar a $2{,}04$~m y el paso sería de $2{,}13$~m: cabrían 6 filas, es decir, 96 paneles. Esa fila de más quedaría a la sombra de la anterior en las primeras y últimas horas de los días de invierno, que es justo lo que el criterio del IDAE quiere evitar. Se elige la propuesta de 80 paneles.

\paso{Paso 17. Producción estimada con PVGIS}
En PVGIS \parencite{pvgis} se introduce la ubicación ($40{,}4^\circ$~N, $3{,}7^\circ$~O), la potencia de 44~kWp, una inclinación de $30^\circ$, orientación sur y las pérdidas por defecto del 14\,\%. La herramienta estima unos \textbf{71\,300~kWh al año}, unos 1\,620~kWh por cada kWp instalado, con un mínimo de unos 4\,300~kWh en noviembre y un máximo de unos 7\,600~kWh en julio.

\paso{Paso 18. Incidencia sobre los paneles}
Con $\beta_p=30^\circ$, la fórmula $\cos\theta=\sin(\beta_p+h)$ (resultado~\ref{res:incidencia}) da al mediodía:

\begin{tablaSA}{colspec={X[1.6,l]X[1,c]X[1,c]X[1,c]}, cells={font=\footnotesize}}
  Fecha & $h$ & $\cos\theta=\sin(30^\circ+h)$ & $\theta$ \\
  Solsticio de invierno & $26{,}15^\circ$ & 0,83 & $33{,}85^\circ$ \\
  Equinoccio & $49{,}60^\circ$ & 0,98 & $10{,}40^\circ$ \\
  Solsticio de verano & $73{,}05^\circ$ & 0,97 & $13{,}05^\circ$ \\
\end{tablaSA}

La incidencia sería perfecta con $\beta_p+h=90^\circ$. Una inclinación de $30^\circ$ favorece la primavera y el verano, cuando más produce la instalación, y además reduce la altura de las filas y, con ella, la separación necesaria.

% ======================================================================
\section{Reflexionar sobre el modelo (S16)}

\paso{Paso 19. Aportación a los ODS}
La propuesta añade unos 71\,300~kWh anuales de energía renovable producida en el propio centro (ODS~7, meta~7.2) y aprovecha una cubierta ya construida en lugar de ocupar el patio, que sigue disponible para la comunidad (ODS~11).

\paso{Paso 20. Límites del modelo}
\begin{itemize}
  \item El plano sur-norte describe el mediodía; a otras horas la sombra se tuerce (resultado~\ref{res:avance}) y por eso se usa el criterio del IDAE.
  \item Se supone el edificio orientado de este a oeste y el terreno llano.
  \item Puede haber obstáculos no medidos: chimeneas, antenas, árboles o castilletes en la cubierta.
  \item No se consideran la nubosidad real de cada día, el coste, la resistencia de la cubierta ni los permisos.
  \item Las alturas tienen un error de medida que se ha acotado, pero no eliminado.
\end{itemize}

% ======================================================================
\section{Cerrar el libro y defender (S17 a S20)}\label{sec:cierre}

\paso{Paso 21. Registro de cambios}
Ejemplo de entradas tras la coevaluación de S16:

\begin{tablaSA}{colspec={X[1.3,l]X[1.2,l]X[1.3,l]}, cells={font=\footnotesize}}
  Crítica recibida & Qué se modifica & Por qué \\
  «No se entiende por qué la franja no sirve» & Se añade el cálculo con el peor y el mejor caso del error & La conclusión debe ser robusta frente al error de medida \\
  «Las filas parecen demasiado separadas» & Se incluye la comparación con el criterio del mediodía (96 paneles) & Hay que justificar por qué se renuncia a una fila \\
  «Falta la unidad en la tabla de alturas» & Se revisan unidades en todas las tablas & Notación correcta (criterio 8.1) \\
\end{tablaSA}

\paso{Paso 22. Fuentes y autoría}
El libro cita la fórmula de la declinación \parencite{pveducationdecl}, la de la altura del sol \parencite{pveducation}, el criterio de separación entre filas \parencite{idae}, la estimación de producción \parencite{pvgis} y, si se usan, los datos de la calculadora de la NOAA \parencite{noaa}.

\paso{Paso 23. Estructura de la defensa (5 minutos)}
Problema y zonas estudiadas (1~min); medición y error (1~min); modelo y validación con la varilla (1~min); propuesta, 80 paneles y producción (1{,}5~min); límites y conclusión (0{,}5~min). Cada miembro del equipo presenta una parte.

% ======================================================================
\part*{Demostraciones}
\addcontentsline{toc}{part}{Demostraciones}

\section{Por qué funcionan las fórmulas}\label{sec:demos}

Los resultados siguen un orden lógico: cada uno usa solo los anteriores. Se trabaja con ángulos en grados y se usan las razones de un ángulo cualquiera definidas en la circunferencia goniométrica.

\subsection{Trigonometría del triángulo}

\resultado{Las razones no dependen del triángulo}{Si dos triángulos rectángulos tienen un ángulo agudo igual a $\alpha$, el cociente entre dos lados correspondientes es el mismo en ambos. Por eso $\sin\alpha$, $\cos\alpha$ y $\tan\alpha$ dependen solo de $\alpha$.}\label{res:razones}
\begin{dem}
Los dos triángulos tienen un ángulo recto y un ángulo $\alpha$, así que el tercer ángulo es $90^\circ-\alpha$ en ambos. Con los tres ángulos iguales son semejantes, y sus lados son proporcionales: $a'=\lambda a$, $b'=\lambda b$ y $c'=\lambda c$. Por tanto $a'/c'=\lambda a/(\lambda c)=a/c$, y lo mismo ocurre con los otros cocientes.
\end{dem}

\resultado{Altura con un punto de medida}{Con las notaciones de la figura~\ref{fig:alturas}a, $H=D\tan\alpha+a$.}\label{res:unpunto}
\begin{dem}
El triángulo formado por el ojo, el punto $F$ de la fachada a la altura del ojo y lo alto $T$ es rectángulo en $F$. El cateto adyacente a $\alpha$ mide $D$ y el opuesto, $x=FT$. Por definición de tangente, $\tan\alpha=x/D$, así que $x=D\tan\alpha$. Como $F$ está a la altura $a$ del suelo, $H=x+a$.
\end{dem}

\resultado{Teorema del coseno}{En un triángulo de lados $a$, $b$, $c$, con $C$ el ángulo opuesto a $c$, se cumple $c^2=a^2+b^2-2ab\cos C$.}\label{res:coseno}
\begin{dem}
Se colocan los vértices en $C=(0,0)$, $B=(a,0)$ y, por la definición de las razones en la circunferencia goniométrica, $A=(b\cos C,\ b\sin C)$. Esto vale también si $C$ es obtuso. Entonces
\[
  c^2=|AB|^2=(b\cos C-a)^2+(b\sin C)^2=b^2(\cos^2 C+\sin^2 C)-2ab\cos C+a^2=a^2+b^2-2ab\cos C.
\]
\end{dem}

\resultado{Teorema del seno}{En cualquier triángulo, $\dfrac{a}{\sin A}=\dfrac{b}{\sin B}=\dfrac{c}{\sin C}$.}\label{res:seno}
\begin{dem}
Sea $h_c$ la altura desde $C$ sobre la recta $AB$. En el triángulo rectángulo que forma con el lado $b$, $h_c=b\sin A$. En el que forma con el lado $a$, $h_c=a\sin B$. Si $B$ es obtuso, el pie de la altura cae fuera del lado $AB$ y el ángulo de ese triángulo es $180^\circ-B$, pero $\sin(180^\circ-B)=\sin B$. Igualando, $b\sin A=a\sin B$, es decir, $a/\sin A=b/\sin B$. Con la altura desde $A$ se obtiene la igualdad con $c/\sin C$.
\end{dem}

\subsection{Vectores y fórmulas de adición}

\resultado{Expresión geométrica del producto escalar}{Si $\vec u$ y $\vec v$ forman un ángulo $\theta$, entonces $u_1v_1+u_2v_2=|\vec u|\,|\vec v|\cos\theta$.}\label{res:escalar}
\begin{dem}
Los vectores $\vec u$, $\vec v$ y $\vec u-\vec v$ forman un triángulo con ángulo $\theta$ entre los dos primeros. Por el teorema del coseno, $|\vec u-\vec v|^2=|\vec u|^2+|\vec v|^2-2|\vec u||\vec v|\cos\theta$. En coordenadas, $|\vec u-\vec v|^2=(u_1-v_1)^2+(u_2-v_2)^2=|\vec u|^2+|\vec v|^2-2(u_1v_1+u_2v_2)$. Igualando las dos expresiones se obtiene el resultado.
\end{dem}
En particular, dos vectores son perpendiculares si y solo si $\vec u\cdot\vec v=0$, y la proyección ortogonal de $\vec v$ sobre un vector unitario $\vec e$ mide $\vec v\cdot\vec e=|\vec v|\cos\theta$.

\resultado{Coseno de la diferencia y seno de la suma}{Para ángulos cualesquiera $a$ y $b$:
\[
  \cos(a-b)=\cos a\cos b+\sin a\sin b,\qquad \sin(a+b)=\sin a\cos b+\cos a\sin b,\qquad \sin(a-b)=\sin a\cos b-\cos a\sin b.
\]}\label{res:adicion}
\begin{dem}
Los vectores unitarios $\vec u=(\cos a,\sin a)$ y $\vec v=(\cos b,\sin b)$ forman un ángulo $a-b$ (o su opuesto, que tiene el mismo coseno). Por el resultado~\ref{res:escalar}, $\cos(a-b)=\vec u\cdot\vec v=\cos a\cos b+\sin a\sin b$. Para el seno de la suma se usa que $\sin x=\cos(90^\circ-x)$:
\[
  \sin(a+b)=\cos\big((90^\circ-a)-b\big)=\cos(90^\circ-a)\cos b+\sin(90^\circ-a)\sin b=\sin a\cos b+\cos a\sin b.
\]
Cambiando $b$ por $-b$, y como $\cos(-b)=\cos b$ y $\sin(-b)=-\sin b$, se obtiene el seno de la diferencia.
\end{dem}

\subsection{Alturas inaccesibles}

\resultado{Altura con dos puntos de medida}{Con las notaciones de la figura~\ref{fig:alturas}b y $\alpha<\beta$,
\[
  H=\frac{d\,\sin\alpha\,\sin\beta}{\sin(\beta-\alpha)}+a=\frac{d\tan\alpha\tan\beta}{\tan\beta-\tan\alpha}+a.
\]}\label{res:dospuntos}
\begin{dem}
En el triángulo $ABT$, el ángulo en $A$ es $\alpha$. El ángulo en $B$ es $180^\circ-\beta$, porque $\beta$ es su ángulo exterior. Como los tres suman $180^\circ$, el ángulo en $T$ es $\beta-\alpha$. Por el teorema del seno, $BT/\sin\alpha=d/\sin(\beta-\alpha)$, así que $BT=d\sin\alpha/\sin(\beta-\alpha)$. En el triángulo rectángulo $BFT$, $BT$ es la hipotenusa y $x=FT$ el cateto opuesto a $\beta$, así que $x=BT\sin\beta$. Sumando la altura del ojo se obtiene la primera fórmula.

Para la segunda, sea $D=BF$. En los triángulos rectángulos $BFT$ y $AFT$, $x=D\tan\beta$ y $x=(D+d)\tan\alpha$. Igualando, $D=d\tan\alpha/(\tan\beta-\tan\alpha)$ y $x=d\tan\alpha\tan\beta/(\tan\beta-\tan\alpha)$. Ambas expresiones coinciden, porque por el resultado~\ref{res:adicion}
\[
  \tan\beta-\tan\alpha=\frac{\sin\beta\cos\alpha-\cos\beta\sin\alpha}{\cos\alpha\cos\beta}=\frac{\sin(\beta-\alpha)}{\cos\alpha\cos\beta}.
\]
\end{dem}

\resultado{Sensibilidad al error (ampliación con derivadas)}{Un error pequeño $\Delta\alpha$, en radianes, cambia la altura sobre el ojo aproximadamente en
\[
  \Delta x\approx\frac{D}{\cos^2\alpha}\,\Delta\alpha\ \text{(un punto)},\qquad \Delta x\approx x\,\big(\cot\alpha+\cot(\beta-\alpha)\big)\,\Delta\alpha\ \text{(dos puntos)}.
\]}\label{res:error}
\begin{dem}
Con un punto, $x=D\tan\alpha$ y $x'(\alpha)=D/\cos^2\alpha$. Con dos, $\ln x=\ln d+\ln\sin\alpha+\ln\sin\beta-\ln\sin(\beta-\alpha)$; derivando respecto de $\alpha$,
\[
  \frac{x'}{x}=\frac{\cos\alpha}{\sin\alpha}+\frac{\cos(\beta-\alpha)}{\sin(\beta-\alpha)}=\cot\alpha+\cot(\beta-\alpha).
\]
En ambos casos $\Delta x\approx x'\,\Delta\alpha$.
\end{dem}
Con los datos del ejemplo y $\Delta\alpha=1^\circ\approx 0{,}0175$~rad: con un punto, $\Delta x\approx 24{,}8\cdot 0{,}0175/\cos^2 25^\circ\approx 0{,}53$~m; con dos, $\cot 25^\circ+\cot 13^\circ\approx 6{,}48$ y $\Delta x\approx 11{,}57\cdot 6{,}48\cdot 0{,}0175\approx 1{,}31$~m, en línea con la tabla del paso~9. El término $\cot(\beta-\alpha)$ crece sin límite cuando $\beta-\alpha$ se acerca a $0^\circ$: por eso conviene separar los puntos de medida.

\subsection{El sol y las sombras}

\begin{figure}[H]
\centering
\begin{tikzpicture}[>=Stealth,font=\footnotesize,scale=1.15]
  \draw[thick] (0,0) circle (2);
  \draw[dashed] (-2.6,0) node[left]{ecuador} -- (2.6,0);
  \draw[dotted] (0,-2.4) -- (0,2.5) node[above]{eje (N)};
  \coordinate (P) at (40:2); \coordinate (Q) at (15:2);
  \draw (0,0) -- (P) (0,0) -- (Q);
  \fill (P) circle (1.5pt) node[left=2pt]{$P$}; \fill (Q) circle (1.5pt) node[below=2pt]{$S$};
  \draw (0.7,0) arc (0:40:0.7); \node at (20:0.95) {$\varphi$};
  \draw (1.25,0) arc (0:15:1.25); \node at (7:1.5) {$\delta$};
  \draw[->] (P) -- (2.604,2.185) node[above]{cenit};
  \draw[thick] (0.761,2.205) -- (2.303,0.367) node[right]{horizonte};
  \draw[->,urjc,thick] (4.816,2.166) -- (P); \draw[->,urjc,thick] (4.540,1.217) -- (Q);
  \node[urjc,align=center] at (5.2,1.6) {rayos del sol\\(paralelos)};
  \draw (P) ++(-50:0.55) arc (-50:15:0.55); \path (P) ++(-17:0.78) node{$h$};
  \draw (P) ++(15:1.0) arc (15:40:1.0); \path (P) ++(26:1.45) node{$\varphi-\delta$};
\end{tikzpicture}
\caption{Corte de la Tierra por el meridiano del centro al mediodía solar. $P$ es el centro, a latitud $\varphi$; $S$ es el punto donde el sol está en el cenit, a latitud $\delta$.}\label{fig:tierra}
\end{figure}

\resultado{Altura del sol al mediodía}{Si la Tierra es esférica y los rayos del sol llegan paralelos, al mediodía solar en un lugar de latitud $\varphi\ge\delta$ el sol está a una altura $h=90^\circ-\varphi+\delta$.}\label{res:altura}
\begin{dem}
Al mediodía solar el sol está en el plano del meridiano del lugar, que es el de la figura~\ref{fig:tierra}. En el punto $S$, de latitud $\delta$, el sol está en el cenit, así que los rayos son paralelos a la recta $OS$. En $P$ el cenit sigue la recta $OP$. El ángulo entre $OP$ y $OS$ es $\varphi-\delta$, porque las dos latitudes se miden desde el ecuador en el mismo plano. Como los rayos que llegan a $P$ son paralelos a $OS$, el ángulo entre el cenit de $P$ y la dirección del sol también es $\varphi-\delta$. El horizonte es perpendicular al cenit, así que $h=90^\circ-(\varphi-\delta)$.
\end{dem}
En Madrid siempre se cumple $\varphi\ge\delta$, porque $\delta$ nunca supera $23{,}45^\circ$. Una consecuencia útil para el proyecto eTwinning: el mismo día y al mediodía, la diferencia de alturas del sol entre dos centros es igual a la diferencia de sus latitudes, que es la idea con la que Eratóstenes midió la Tierra. La fórmula de $\delta$ en función del día es un modelo astronómico aproximado que se toma de PVEducation \parencite{pveducationdecl}; no se demuestra aquí.

\resultado{El rayo como recta y los puntos en sombra}{En el plano sur-norte, con el obstáculo ocupando $x\le 0$ hasta la altura $H$, el rayo que roza $T=(0,H)$ es la recta $y=H-x\tan h$, y un punto $P=(x_P,y_P)$ con $x_P>0$ está en sombra si y solo si $y_P<H-x_P\tan h$.}\label{res:rayo}
\begin{dem}
La luz avanza hacia el norte y hacia abajo, en la dirección de $\vec v=(\cos h,-\sin h)$, así que el rayo es $T+t\vec v=(t\cos h,\ H-t\sin h)$ con $t\ge 0$. Despejando $t=x/\cos h$ queda $y=H-x\tan h$.

Desde $P$, la dirección hacia el sol es $\vec s=(-\cos h,\sin h)$. La semirrecta $P+t\vec s$ llega a la fachada $x=0$ cuando $t=x_P/\cos h$, a la altura $y^*=y_P+x_P\tan h$. El obstáculo tapa el sol si y solo si esa semirrecta choca con él, es decir, si $y^*<H$, que equivale a $y_P<H-x_P\tan h$.
\end{dem}

\resultado{Longitud de la sombra}{Un obstáculo de altura $H$ proyecta sobre el suelo horizontal una sombra de longitud $L=H/\tan h$, y el coeficiente $k=L/H=1/\tan h$ solo depende de la altura del sol.}\label{res:sombra}
\begin{dem}
Un punto del suelo, $y_P=0$, está en sombra si y solo si $0<H-x_P\tan h$, es decir, si $x_P<H/\tan h$ (resultado~\ref{res:rayo}). La sombra va, por tanto, de $x=0$ a $x=H/\tan h$.
\end{dem}
Como la tangente es creciente entre $0^\circ$ y $90^\circ$, cuanto más bajo está el sol, mayor es $k$ y más larga es la sombra. Por eso el solsticio de invierno es el caso que decide la separación.

\resultado{Avance de la sombra hacia el norte}{Si el sol forma un ángulo $A$ con el sur, una sombra de longitud $L$ avanza hacia el norte $L\cos A$.}\label{res:avance}
\begin{dem}
Vista en planta, la sombra apunta en sentido contrario al sol, así que forma un ángulo $A$ con el norte. Si $\vec w$ es el vector unitario en la dirección de la sombra y $\vec e$ el vector unitario hacia el norte, el avance es la proyección ortogonal de $L\vec w$ sobre $\vec e$: $L\,\vec w\cdot\vec e=L\cos A$ (resultado~\ref{res:escalar}).
\end{dem}
La altura y el ángulo a otras horas se toman de la calculadora de la NOAA. Como comprobación optativa, la trigonometría esférica, que no se demuestra aquí, da $\sin h=\sin\varphi\sin\delta+\cos\varphi\cos\delta\cos\omega$ y $\sin A=\cos\delta\sin\omega/\cos h$, donde $\omega$ es el ángulo horario ($15^\circ$ por hora desde el mediodía). Dos horas después del mediodía del solsticio de invierno, $\omega=30^\circ$ y resultan $h\approx 20{,}31^\circ$ y $A\approx 29{,}3^\circ$, los mismos valores de la situación de aprendizaje.

\resultado{Incidencia sobre un panel}{Un panel orientado al sur con inclinación $\beta_p$ recibe el sol de altura $h$ con un ángulo $\theta$ que cumple $\cos\theta=\sin(\beta_p+h)$. La incidencia es perpendicular si y solo si $\beta_p+h=90^\circ$.}\label{res:incidencia}
\begin{dem}
En el plano sur-norte, con $x$ hacia el norte, la superficie del panel sube hacia el norte en la dirección $\vec p=(\cos\beta_p,\sin\beta_p)$. El vector $\vec n=(-\sin\beta_p,\cos\beta_p)$ es unitario, perpendicular a $\vec p$ porque $\vec n\cdot\vec p=0$, y apunta hacia arriba y hacia el sur: es la normal. El vector unitario hacia el sol es $\vec s=(-\cos h,\sin h)$. Por los resultados~\ref{res:escalar} y~\ref{res:adicion},
\[
  \cos\theta=\vec n\cdot\vec s=\sin\beta_p\cos h+\cos\beta_p\sin h=\sin(\beta_p+h).
\]
El coseno vale 1, es decir, $\theta=0^\circ$, si y solo si $\sin(\beta_p+h)=1$, lo que con $0^\circ<\beta_p+h<180^\circ$ ocurre solo si $\beta_p+h=90^\circ$.
\end{dem}

\resultado{Paso entre filas}{Si cada fila tiene una longitud $l$ en la pendiente y una inclinación $\beta_p$, y la separación se diseña con un sol de altura $h_d$, el paso entre el comienzo de dos filas consecutivas es $p=l\cos\beta_p+k\,l\sin\beta_p$, con $k=1/\tan h_d$.}\label{res:filas}
\begin{dem}
La fila delantera alcanza una altura $h_o=l\sin\beta_p$ y ocupa en planta $l\cos\beta_p$. Para que no haga sombra a la siguiente, el pie de esta debe quedar donde termina la sombra de la delantera con el sol de diseño, es decir, a $h_o/\tan h_d=k\,h_o$ de su borde trasero (resultado~\ref{res:sombra}). Sumando la ocupación en planta y esa separación se obtiene $p$ (figura~\ref{fig:filas}).
\end{dem}
El pliego del IDAE fija $h_d=61^\circ-\varphi$ \parencite{idae}; es un criterio técnico, no un resultado que se demuestre. A $40{,}4^\circ$ da $h_d=20{,}6^\circ$, unos $5{,}55^\circ$ por debajo de la altura del mediodía del solsticio de invierno, $66{,}55^\circ-\varphi$. Por el resultado~\ref{res:sombra}, un sol más bajo da un $k$ mayor ($2{,}66$ frente a $2{,}04$): ese es el margen que protege las filas fuera del mediodía.

% ======================================================================
\printbibliography[heading=bibintoc,title={Bibliografía}]

\section*{Declaración de uso de inteligencia artificial}
\addcontentsline{toc}{section}{Declaración de uso de inteligencia artificial}
Para la elaboración de este documento se ha utilizado una herramienta de inteligencia artificial generativa (Claude, de Anthropic) para redactar el solucionario y las demostraciones, generar el código \LaTeX{} y las figuras, y calcular y comprobar los valores numéricos. El alumno ha revisado y validado el contenido y asume la plena responsabilidad sobre el documento.

\end{document}
